\documentclass[12pt,a4paper]{article}
\usepackage[utf8]{inputenc}
\usepackage[T1]{fontenc}
\usepackage[french]{babel}
\usepackage{geometry}
\usepackage{titlesec}
\usepackage{parskip}
\usepackage{hyperref}
\usepackage{xcolor}
\usepackage{enumitem}
\usepackage{fancyhdr}
\usepackage{booktabs}
\usepackage{array}

\geometry{margin=2.5cm}

\definecolor{codeblue}{RGB}{0,70,150}
\definecolor{codered}{RGB}{180,0,0}
\definecolor{codegray}{RGB}{100,100,100}

\pagestyle{fancy}
\fancyhf{}
\rhead{Application de statistiques sportives}
\lhead{Module \texttt{src/Parsers}}
\cfoot{\thepage}

\titleformat{\section}{\large\bfseries\color{codeblue}}{}{0em}{}[\titlerule]
\titleformat{\subsection}{\normalsize\bfseries\color{codered}}{}{0em}{}
\titleformat{\subsubsection}{\normalsize\itshape}{}{0em}{}

\newcommand{\fonction}[1]{\textbf{\texttt{#1}}}
\newcommand{\module}[1]{\textit{\texttt{#1}}}

\begin{document}

\begin{titlepage}
    \centering
    \vspace*{3cm}
    {\huge\bfseries Documentation technique\\[0.5em]}
    {\Large\bfseries Module \texttt{src/Parsers}\\[2em]}
    {\large Application de statistiques sportives\\[1em]}
    {\normalsize Projet 1A — 2026\\[3em]}
    \hrule
    \vspace{1em}
    {\small Ce document décrit l'ensemble des classes et fonctions du dossier
    \texttt{src/Parsers}, responsable du chargement et du nettoyage des données
    sportives depuis les fichiers CSV.}
\end{titlepage}

\tableofcontents
\newpage

% ============================================================
\section{Architecture générale}

Le dossier \texttt{src/Parsers} repose sur un pattern d'héritage unique~:
une classe mère \texttt{BaseLoader} centralise toute la mécanique commune,
et chaque sport hérite d'elle en ajoutant uniquement ce qui lui est spécifique.

\begin{center}
\begin{tabular}{l l}
\toprule
\textbf{Classe} & \textbf{Sport couvert} \\
\midrule
\texttt{BaseLoader} & Classe abstraite — socle commun \\
\texttt{BadmintonLoader} & BWF World Tour \\
\texttt{BasketballLoader} & NBA 2022-2023 \\
\texttt{ChessLoader} & FIDE (Échecs) \\
\texttt{Cs2Loader} & Counter-Strike 2 \\
\texttt{FootballLoader} & Ligues européennes 2008-2016 \\
\texttt{FootballChampionsLeagueLoader} & UEFA Champions League \\
\texttt{LolLoader} & League of Legends EMEA 2025 \\
\texttt{Starcraft2Loader} & StarCraft II \\
\texttt{TennisLoader} & ATP et WTA 2024 \\
\texttt{VolleyballLoader} & JO Paris 2024 \\
\bottomrule
\end{tabular}
\end{center}

Le fichier \texttt{\_\_init\_\_.py} expose les dix loaders comme imports directs
du package, permettant d'écrire \texttt{from src.Parsers import BasketballLoader}
plutôt que le chemin complet.

% ============================================================
\section{\texttt{BaseLoader} — Classe mère abstraite}

\textbf{Fichier~:} \texttt{src/Parsers/BaseLoader.py}

\texttt{BaseLoader} hérite de \texttt{ABC} (classe abstraite Python), ce qui
interdit de l'instancier directement. Toutes les sous-classes doivent déclarer
l'attribut \texttt{SPORT\_FOLDER}.

\subsection{Attributs de classe}

\begin{itemize}[leftmargin=1.5cm]
    \item \fonction{\_DATA\_ROOT} — Chemin absolu vers le dossier \texttt{data/},
    calculé dynamiquement depuis l'emplacement du fichier via
    \texttt{Path(\_\_file\_\_).parent.parent.parent / "data"}.
    Ce calcul garantit le fonctionnement quel que soit le répertoire de lancement.
    \item \fonction{SPORT\_FOLDER} — Chaîne vide par défaut.
    Doit être redéfinie dans chaque sous-classe (ex~: \texttt{"Basketball"},
    \texttt{"tennis"}).
\end{itemize}

\subsection{Propriété \texttt{root}}

Retourne le chemin complet vers le dossier de données du sport~:
\texttt{\_DATA\_ROOT / SPORT\_FOLDER}. Déclarée comme \texttt{@property},
elle s'utilise sans parenthèses (\texttt{self.root}).

\subsection{Méthode \texttt{\_load\_csv(filename, dtype, date\_cols)}}

Méthode privée centrale du loader. Elle~:
\begin{enumerate}
    \item Lit le fichier CSV via \texttt{pd.read\_csv()} en appliquant les types
    fournis dans \texttt{dtype}.
    \item Convertit automatiquement chaque colonne listée dans \texttt{date\_cols}
    via \texttt{pd.to\_datetime(..., errors="coerce")}~: une date malformée
    devient \texttt{NaT} sans lever d'exception.
\end{enumerate}

% ============================================================
\section{\texttt{BadmintonLoader}}

\textbf{Fichier~:} \texttt{src/Parsers/BadmintonLoader.py}\\
\textbf{Données~:} BWF World Tour — 88 joueurs, 226 matchs.\\
\textbf{SPORT\_FOLDER~:} \texttt{"badminton"}

\subsection{\texttt{load\_players()}}
Charge \texttt{player.csv}. Colonnes retournées~: \texttt{name}, \texttt{country},
\texttt{continent}. Aucune transformation appliquée.

\subsection{\texttt{load\_matches()}}
Charge \texttt{match.csv} avec conversion de la colonne \texttt{date} en datetime.
Colonnes~: \texttt{tournament}, \texttt{city}, \texttt{country}, \texttt{date},
\texttt{tournament\_type}, \texttt{round}, \texttt{player\_1}, \texttt{player\_2},
\texttt{winner}, \texttt{game\_1\_score}, \texttt{game\_2\_score},
\texttt{game\_3\_score}.

\subsection{\texttt{load\_all()}}
Appelle \texttt{load\_players()} et \texttt{load\_matches()} et retourne le
tuple \texttt{(players, matches)}.

% ============================================================
\section{\texttt{BasketballLoader}}

\textbf{Fichier~:} \texttt{src/Parsers/BasketballLoader.py}\\
\textbf{Données~:} NBA saison 2022-2023 — 431 joueurs, 30 équipes.\\
\textbf{SPORT\_FOLDER~:} \texttt{"Basketball"}

\subsection{\texttt{load\_players()}}
Charge \texttt{player.csv} et ajoute deux colonnes calculées~:
\begin{itemize}
    \item \texttt{full\_name} — concaténation \texttt{first\_name + " " + last\_name}.
    \item \texttt{height\_cm} — conversion du format américain pieds-pouces
    (\texttt{"6-8"}) en centimètres via \texttt{\_height\_to\_cm()}.
\end{itemize}

\subsection{\texttt{load\_teams()}}
Charge \texttt{team.csv} avec typage forcé \texttt{id → int}. Colonnes~:
\texttt{id}, \texttt{full\_name}, \texttt{abbreviation}, \texttt{nickname},
\texttt{city}, \texttt{state}.

\subsection{\texttt{load\_matches()}}
Charge \texttt{game.csv} avec statistiques complètes par équipe~: points,
rebonds, passes décisives, tirs, lancers francs, etc.

\subsection{\texttt{load\_all()}}
Retourne le tuple \texttt{(players, teams, matches)}.

\subsection{\texttt{get\_team\_roster(players, teams, team\_name)}}
Méthode utilitaire. Recherche une équipe par nom complet, abréviation ou surnom
(insensible à la casse), puis filtre les joueurs correspondants.
Lève une \texttt{ValueError} si aucune équipe ne correspond.

\subsection{\texttt{\_height\_to\_cm(height\_str)}}
Méthode statique privée. Convertit le format pieds-pouces \texttt{"6-8"} en
centimètres~: \textit{(pieds × 12 + pouces) × 2,54}. Retourne \texttt{None}
si la valeur est manquante.

% ============================================================
\section{\texttt{ChessLoader}}

\textbf{Fichier~:} \texttt{src/Parsers/ChessLoader.py}\\
\textbf{Données~:} FIDE — 206 joueurs, 256 matchs.\\
\textbf{SPORT\_FOLDER~:} \texttt{"chess"}

\subsection{\texttt{load\_players()}}
Charge \texttt{player.csv} avec typage \texttt{Int64} (entier nullable) pour les
colonnes Elo~: \texttt{rating\_standard}, \texttt{rating\_rapid},
\texttt{rating\_blitz}. Le type nullable gère proprement les joueurs sans
classement dans une catégorie.

\subsection{\texttt{load\_matches()}}
Charge \texttt{match.csv}. Particularité~: la colonne \texttt{score\_player\_1}
peut contenir la valeur textuelle \texttt{"Bye"} (joueur exempt de tour), ce
qui justifie le typage \texttt{Float64} pour certaines colonnes.

\subsection{\texttt{load\_all()}}
Retourne le tuple \texttt{(players, matches)}.

% ============================================================
\section{\texttt{Cs2Loader}}

\textbf{Fichier~:} \texttt{src/Parsers/Cs2Loader.py}\\
\textbf{Données~:} Counter-Strike 2 — 160 joueurs, 32 coachs, 32 équipes,
106 matchs.\\
\textbf{SPORT\_FOLDER~:} \texttt{"counter\_strike\_2"}

\subsection{\texttt{load\_players()}}
Charge \texttt{player.csv} avec conversion de \texttt{birthdate}.
Colonnes~: \texttt{pseudo}, \texttt{name}, \texttt{nationality},
\texttt{birthdate}, \texttt{role}, \texttt{team}.

\subsection{\texttt{load\_coaches()}}
Charge \texttt{coach.csv} avec conversion de \texttt{birthdate}.
Structure identique à \texttt{load\_players()}.

\subsection{\texttt{load\_teams()}}
Charge \texttt{team.csv}. Colonnes~: \texttt{team}, \texttt{team\_abbreviation},
\texttt{location}, \texttt{region}.

\subsection{\texttt{load\_matches()}}
Charge \texttt{match.csv} avec conversion de \texttt{date}. Colonnes~:
\texttt{date}, \texttt{stage}, \texttt{round}, \texttt{best\_of},
\texttt{team\_1}, \texttt{team\_2}, \texttt{score\_team\_1},
\texttt{score\_team\_2}.

\subsection{\texttt{load\_all()}}
Retourne le tuple \texttt{(players, coaches, teams, matches)}.

% ============================================================
\section{\texttt{FootballLoader}}

\textbf{Fichier~:} \texttt{src/Parsers/FootballLoader.py}\\
\textbf{Données~:} Ligues européennes — 25\,979 matchs sur les saisons
2008-2016.\\
\textbf{SPORT\_FOLDER~:} \texttt{"Football"}

C'est le loader avec le plus grand volume de données. Il charge cinq tables
distinctes.

\subsection{\texttt{load\_countries()}}
Charge \texttt{country.csv}. Colonnes~: \texttt{id}, \texttt{name}.

\subsection{\texttt{load\_leagues()}}
Charge \texttt{league.csv}. Colonnes~: \texttt{id}, \texttt{country\_id},
\texttt{name}.

\subsection{\texttt{load\_teams()}}
Charge \texttt{team.csv}. Colonnes~: \texttt{id}, \texttt{team\_api\_id},
\texttt{team\_long\_name}, \texttt{team\_short\_name}.

\subsection{\texttt{load\_players()}}
Charge \texttt{player.csv}. Colonnes~: \texttt{id}, \texttt{player\_api\_id},
\texttt{player\_name}, \texttt{birthday}, \texttt{weight} (kg),
\texttt{height} (cm).

\subsection{\texttt{load\_matches()}}
Charge \texttt{match.csv} avec conversion de \texttt{date}. Contient les
compositions d'équipe complètes (\texttt{home\_player\_1} à
\texttt{home\_player\_11}).

\subsection{\texttt{load\_all()}}
Retourne uniquement \texttt{(countries, leagues, matches)}. Les tables joueurs
et équipes sont disponibles via leurs méthodes dédiées mais non incluses dans
\texttt{load\_all()}.

% ============================================================
\section{\texttt{FootballChampionsLeagueLoader}}

\textbf{Fichier~:} \texttt{src/Parsers/FootballChampionsLeagueLoader.py}\\
\textbf{Données~:} UEFA Champions League — 751 joueurs, 32 équipes,
125 matchs.\\
\textbf{SPORT\_FOLDER~:} \texttt{"football\_champions\_league"}

\subsection{\texttt{load\_players()}}
Charge \texttt{player.csv} et corrige un problème dans le fichier source~:
deux colonnes portent le même nom \texttt{match\_played}. Le loader les renomme
en \texttt{match\_played\_field} et \texttt{match\_played} via
\texttt{rename()}. Des fautes de frappe d'origine
(\texttt{balls\_recoverd}, \texttt{cross\_complted}) sont conservées
volontairement pour rester fidèle au fichier source.

\subsection{\texttt{load\_teams()}}
Charge \texttt{team.csv}. Colonnes~: \texttt{full\_name}, \texttt{short\_name},
\texttt{year\_founded}, \texttt{country}, \texttt{league}, \texttt{city}.

\subsection{\texttt{load\_matches()}}
Charge \texttt{match.csv} avec conversion de \texttt{date}. Colonnes~:
\texttt{date}, \texttt{phase}, \texttt{round}, \texttt{matchday},
\texttt{group}, \texttt{team\_home}, \texttt{team\_away}, scores.

\subsection{\texttt{load\_all()}}
Retourne le tuple \texttt{(players, teams, matches)}.

% ============================================================
\section{\texttt{LolLoader}}

\textbf{Fichier~:} \texttt{src/Parsers/LolLoader.py}\\
\textbf{Données~:} League of Legends EMEA 2025 — 50 joueurs, 25 coachs,
10 équipes, 45 matchs.\\
\textbf{SPORT\_FOLDER~:} \texttt{"Lol"}

\subsection{\texttt{load\_players()}}
Charge \texttt{player.csv} avec conversion de \texttt{birthdate}. Colonnes~:
\texttt{pseudo}, \texttt{name}, \texttt{country\_of\_birth}, \texttt{birthdate},
\texttt{role}, \texttt{team}.

\subsection{\texttt{load\_coaches()}}
Charge \texttt{coach.csv} avec conversion de \texttt{birthdate}.
Structure identique à \texttt{load\_players()}.

\subsection{\texttt{load\_teams()}}
Charge \texttt{team.csv}. Colonnes~: \texttt{team}, \texttt{team\_abbreviation},
\texttt{location}, \texttt{region}.

\subsection{\texttt{load\_matches()}}
Charge \texttt{match.csv} et ajoute une colonne calculée \texttt{duration\_s}
(durée de la partie en secondes), convertie depuis le format texte
\texttt{HH:MM:SS} via \texttt{pd.to\_timedelta().dt.total\_seconds()}.

\subsection{\texttt{load\_all()}}
Retourne le tuple \texttt{(players, coaches, teams, matches)}.

% ============================================================
\section{\texttt{Starcraft2Loader}}

\textbf{Fichier~:} \texttt{src/Parsers/Starcraft2Loader.py}\\
\textbf{Données~:} StarCraft II — 32 joueurs, 67 matchs.\\
\textbf{SPORT\_FOLDER~:} \texttt{"starcraft\_2"}

\subsection{\texttt{load\_players()}}
Charge \texttt{player.csv} avec conversion de \texttt{birthdate}. Colonnes~:
\texttt{pseudo}, \texttt{name}, \texttt{nationality}, \texttt{birthdate},
\texttt{race} (Terran / Zerg / Protoss), \texttt{team}.

\subsection{\texttt{load\_matches()}}
Charge \texttt{match.csv} avec conversion de \texttt{date}. Colonnes~:
\texttt{date}, \texttt{round}, \texttt{group}, \texttt{best\_of},
\texttt{player\_1}, \texttt{player\_2}, \texttt{score\_player\_1},
\texttt{score\_player\_2}.

\subsection{\texttt{load\_all()}}
Retourne le tuple \texttt{(players, matches)}.

% ============================================================
\section{\texttt{TennisLoader}}

\textbf{Fichier~:} \texttt{src/Parsers/TennisLoader.py}\\
\textbf{Données~:} ATP et WTA 2024 — 443 joueurs ATP, 335 joueuses WTA,
environ 5\,700 matchs au total.\\
\textbf{SPORT\_FOLDER~:} \texttt{"tennis"}

Seul loader à gérer deux circuits en parallèle (ATP/WTA), ce qui double le
nombre de méthodes.

\subsection{\texttt{load\_atp\_players()} / \texttt{load\_wta\_players()}}
Chargent les joueurs/joueuses, convertissent \texttt{dob} (format
\texttt{YYYYMMDD} stocké comme float) via \texttt{\_parse\_dob()} et ajoutent
\texttt{full\_name}.

\subsection{\texttt{load\_atp\_matches()} / \texttt{load\_wta\_matches()}}
Chargent les matchs et convertissent \texttt{tourney\_date} (format
\texttt{YYYYMMDD} stocké comme int) via \texttt{\_parse\_tourney\_date()}.

\subsection{\texttt{load\_all()}}
Retourne le tuple \texttt{(atp\_players, wta\_players, atp\_matches,
wta\_matches)}.

\subsection{\texttt{\_parse\_dob(series)}}
Méthode statique privée. Convertit un float \texttt{19860504.0} en
\texttt{datetime(1986, 5, 4)}. Utilise une chaîne
\texttt{dropna().astype(int).astype(str)} pour éviter de planter sur les
valeurs manquantes.

\subsection{\texttt{\_parse\_tourney\_date(series)}}
Méthode statique privée. Même logique que \texttt{\_parse\_dob} mais depuis
un entier directement.

% ============================================================
\section{\texttt{VolleyballLoader}}

\textbf{Fichier~:} \texttt{src/Parsers/VolleyballLoader.py}\\
\textbf{Données~:} JO Paris 2024 — 156 joueurs H, 155 joueuses F, 39 coachs H,
42 coachs F, 26 matchs H, 26 matchs F.\\
\textbf{SPORT\_FOLDER~:} \texttt{"volleyball"}

Loader le plus volumineux avec sept tables et une distinction systématique
hommes/femmes.

\subsection{\texttt{load\_countries()}}
Charge la table des 224 codes IOC. Colonnes~: \texttt{code}, \texttt{country},
\texttt{country\_long}.

\subsection{\texttt{load\_players\_men()} / \texttt{load\_players\_women()}}
Chargent les joueurs/joueuses et ajoutent une colonne \texttt{gender}
(\texttt{"M"} ou \texttt{"F"}) pour permettre la fusion des deux tables sans
ambiguïté.

\subsection{\texttt{load\_coaches\_men()} / \texttt{load\_coaches\_women()}}
Chargent les coachs masculins et féminins avec conversion de \texttt{birth\_date}.

\subsection{\texttt{load\_matches\_men()}}
Charge \texttt{match\_men.csv} et calcule automatiquement le gagnant en
comparant \texttt{set\_country\_1} et \texttt{set\_country\_2}.

\subsection{\texttt{load\_matches\_women()}}
Même logique que pour les hommes, mais le fichier source utilise des noms de
colonnes différents (\texttt{country\_1} au lieu de \texttt{country\_code\_1}).
Le loader les renomme pour assurer l'uniformité avec \texttt{load\_matches\_men()}.

\subsection{\texttt{load\_all()}}
Retourne les sept tables dans l'ordre~: \texttt{(countries, players\_men,
players\_women, coaches\_men, coaches\_women, matches\_men, matches\_women)}.

% ============================================================
\section{Synthèse des choix architecturaux}

\begin{center}
\begin{tabular}{>{\ttfamily}p{5.5cm} p{8cm}}
\toprule
\textbf{Choix} & \textbf{Bénéfice} \\
\midrule
Héritage de BaseLoader & Zéro duplication pour la lecture CSV et les dates \\
SPORT\_FOLDER comme attribut de classe & Ajouter un sport = créer une sous-classe, pas modifier BaseLoader \\
\_DATA\_ROOT calculé dynamiquement & Fonctionne quel que soit le répertoire de lancement \\
errors="coerce" sur les dates & Données manquantes deviennent NaT sans exception \\
Méthodes statiques pour les parseurs de dates (Tennis) & Réutilisables sans instanciation, logique isolée et testable \\
\bottomrule
\end{tabular}
\end{center}

\end{document}
